% Цвет заголовков разделов (тёмно-красный)
\definecolor{cvred}{RGB}{163,29,42}

% Цвет второстепенного текста (например, тип обучения курсивом)
\definecolor{cvgray}{RGB}{90,90,90}

% Заголовок раздела резюме: жирный красный текст + линия под ним
\newcommand{\cvsection}[1]{%
    \vspace{10pt}%
    {\color{cvred}\bfseries\large #1}\par
    \vspace{2pt}
    {\color{cvred}\hrule height 1pt}%
    \vspace{6pt}%
}

% Строка "заголовок слева — период/тип справа" (используется для образования)
\newcommand{\cventry}[2]{%
    \noindent\textbf{#1}\hfill\textit{#2}\\%
}

% Строка проекта: название, необязательная пометка (например, "личный проект"), технологии справа
\newcommand{\cvproject}[3][]{%
    \noindent\textbf{#2}%
    \ifblank{#1}{}{\ \ {\small\itshape #1}}%
    \hfill\textit{#3}\\%
}

% Второстепенная строка курсивом серого цвета (например, тип обучения)
\newcommand{\cvsubtext}[1]{{\color{cvgray}\itshape #1}}

% Нумерованный список проектов/достижений с компактными отступами
\newenvironment{cvlist}{%
    \begin{enumerate}
        \setlength{\itemsep}{5pt}
        \setlength{\parsep}{0pt}
        \setlength{\topsep}{0pt}
}{%
    \end{enumerate}
    \vspace{-8pt}
}
